\begin{textsample}{NEG Dim 9 – global\_north\_2024\_01 – Score -1.00 – global\_north\_2024\_01/104765248.txt}  \label{ex:f9_neg_194}
Market Extra These 3 flashpoints could spark a wider Middle East war that markets can't ignore There is ' clear and present danger ' of a wider conflict involving Iran, says RBC ' s Helima Croft The war between Israel and Hamas that ignited on Oct. 7 has also inflamed tensions on the country ' s northern border with Lebanon, where the militant group Hezbollah has traded cross-border fire with Israel.
Referenced Symbols Oil traders are, for now, looking past rising tensions in the Red Sea.
But that may not last, a top Wall Street commodities analyst warned in a note Tuesday.
Iran ' s decision to send a warship into the Red Sea only briefly alarmed oil traders Tuesday, with crude prices soon dipping back into negative territory after a strong early surge.
Market participants appear to have widely subscribed to the notion that the Israel-Hamas war will remain confined largely to Gaza and don't present a threat to shipments of Middle East crude supplies, said Helima Croft, head of global commodity strategy at RBC Capital Markets, in @ @ @ @ @ @ @ @ @ @ see a clear and present danger that this will become a wider regional conflict, with three flashpoints posing particular contagion concerns,"she wrote."Iran is at the center of all three threat vectors and the question of whether the Gaza conflict will metastasize is likely to hinge on whether Tehran maintains any appreciable degree of distance from its proxies or becomes a main character, either through its own direct actions or because it is targeted for retaliation for its extensive armed network,"she said.
Here are the three potential flashpoints that bear watching, according to Croft : Houthi attacks on shipping The map below points to major incidents around the Red Sea in recent weeks : Meanwhile, there are indications that the U.S. and U.K. are preparing a larger military operation against the Houthis in Yemen, Croft noted.
If so, a key question is whether Iran will remain a bystander.
Croft said it was also noteworthy that Saudi Arabia has reportedly been urging restraint in response to rising Red @ @ @ @ @ @ @ @ @ @ since a cooling of Saudi-Iran tensions last March.
A more serious U.S./Houthi/Iran confrontation could imperil Saudi Arabia ' s newfound security gains, she said.
Lebanon and Hezbollah Senior Biden administration officials have made frequent trips to Beirut recently, in an effort to avert the opening of a second warfront on Israel ' s northern border pitting the nation ' s armed forced against the militant Hezbollah movement.
Israel and Hezbollah have traded fire since the Oct. 7 start of the Israel-Hamas war but actions have remained relatively contained.
Tensions and fighting have increased in recent weeks, however.
An expansion of the conflict into Lebanon would present a"clear pathway"to a wider regional war given Iran ' s extremely close operational relationship with Hezbollah, Croft said, noting that Tehran ' s ties with the group are much stronger than its alliance with the Houthis.
A more-direct Iranian role in the war would sharply increase the threat to regional energy supplies, Croft added -- not only because of Iran ' s energy assets, but also due to @ @ @ @ @ @ @ @ @ @ ( see map below ).
RBC Capital Markets The strait is a narrow waterway that links the Persian Gulf with the Gulf of Oman and the Arabian Sea.
At its narrowest point, the passage is only 21 miles wide, while the width of the shipping lane in either direction is just two miles and separated by a two-mile buffer zone.
It ' s the world ' s most sensitive oil-transportation choke point."While much of the current market focus is on the Red Sea and the Bab el-Mandeb choke point, the Strait of Hormuz is far more consequential for global energy supplies, with an average of 15 million barrels a day of OPEC+ crude traveling through the waterway in 2023,"Croft said.
Iraq firefights Iran-backed militias have fired hundreds of rockets at bases housing U.S. troops and areas near U.S. personnel in Iraq and Syria, prompting a series of escalating retaliatory strikes by the U.S. military, Croft noted."While the militia attacks have injured multiple U.S. soldiers, to date @ @ @ @ @ @ @ @ @ @."We do not know whether President Biden shares his predecessor ' s kinetic redline over the loss of American life.
However, we believe that if continuing assaults by Iran-backed forces lead to significant American casualties, Washington could find itself drawn back into another round of active fighting in Iraq."Oil futures jumped early Tuesday as traders reacted to Iran ' s decision to deploy the warship to the Red Sea, but the gains proved short lived -- with Brent crude BRN00, +0.25 \%, the global benchmark, and West Texas Intermediate crude CL.1, -1.61\%CL00, -1.61 \% both turning south to begin the new year with losses.
About the Author William Watts is MarketWatch markets editor.
In addition to managing markets coverage, he writes about stocks, bonds, \textbf{currencies} and commodities, including oil.
He also writes about global macro issues and trading strategies.
During his time at MarketWatch, Watts has served in key roles in the Frankfurt, London, New York and Washington, D.C., newsrooms.
Dow Jones @ @ @ @ @ @ @ @ @ @ terms of use.
Historical and current end-of-day data provided by FACTSET.
All quotes are in local exchange time.
Real-time last sale data for U.S. stock quotes reflect trades reported through Nasdaq only.
Intraday data delayed at least 15 minutes or per exchange requirements.

% matched lemmas: currency
\end{textsample}
